\documentclass[aspectratio=169]{beamer}

% --- Preamble ---
\usetheme{default}
\usecolortheme{default}

% Essential packages for physics and math
\usepackage{amsmath}
\usepackage{physics}
\usepackage{graphicx}
\usepackage{booktabs} % For professional tables
\usepackage{hyperref}

% Package for inserting Python code snippets
\usepackage{listings}
\usepackage{xcolor}

\definecolor{codegreen}{rgb}{0,0.6,0}
\definecolor{codegray}{rgb}{0.5,0.5,0.5}
\definecolor{codepurple}{rgb}{0.58,0,0.82}
\definecolor{backcolour}{rgb}{0.95,0.87,0.72}

\lstdefinestyle{mystyle}{
    backgroundcolor=\color{backcolour},   
    commentstyle=\color{codegreen},
    keywordstyle=\color{magenta},
    numberstyle=\tiny\color{codegray},
    stringstyle=\color{codepurple},
    basicstyle=\ttfamily\footnotesize,
    breakatwhitespace=false,         
    breaklines=true,                 
    captionpos=b,                    
    keepspaces=true,                 
    numbers=left,                    
    numbersep=5pt,                  
    showspaces=false,                
    showstringspaces=false,
    showtabs=false,                  
    tabsize=2
}
\lstset{style=mystyle}

% --- Title Information ---
\title[Computational Physics]{Metropolis Monte Carlo}
\subtitle{Portfolio Presentation}
\author{Mikkel}
\institute[Aarhus University]{Department of Physics and Astronomy}
\date{\today}

\begin{document}

% Title Slide
\begin{frame}
    \titlepage
\end{frame}

% % Outline Slide
% \begin{frame}{Outline}
%     \tableofcontents
% \end{frame}

\section{Introduction to Metropolis Monte Carlo}
\begin{frame}{Metropolis Monte Carlo}
    
\end{frame}
\begin{frame}{Why Metropolis Monte Carlo?}
    The Metropolis algorithm bypasses the need to evaluate the intractable partition function $Q$ by relying solely on the ratio of target distributions.
    
    \vspace{0.5cm}
    \textbf{Key Advantages:}
    \begin{itemize}
        \item Evaluates only the energy change $\Delta V = V(x') - V(x)$ between states.
        \item Satisfies detailed balance, ensuring the Markov chain converges to the true canonical Boltzmann distribution.
        \item Automatically captures entropy by allowing thermal fluctuations to explore wide basins.
    \end{itemize}
\end{frame}

\section{Algorithm Implementation}
\begin{frame}[fragile]{The Acceptance Criterion}
    The core logic dictates whether a proposed move is accepted or rejected based on the change in energy and the system temperature.
    
    \vspace{0.3cm}
    \begin{lstlisting}[language=Python]
# Metropolis acceptance criterion
if de <= 0 or np.random.rand() < np.exp(-de / kT):
    P = P_new # Accept the move
    accepted_moves += 1
else:
    # Reject the move, stay at current P
    pass 
    \end{lstlisting}
\end{frame}

\section{Simulation Results}
\begin{frame}{Energy Fluctuations and Thermal Average}
    \begin{columns}
        \begin{column}{0.5\textwidth}
            After an initial equilibration period, the system's energy fluctuates stably around the thermal average. 
            
            \vspace{0.3cm}
            These fluctuations can be used directly to calculate macroscopic properties like heat capacity:
            $$C_{V,pot} = \frac{\langle V(x)^2 \rangle - \langle V(x) \rangle^2}{k_B T^2}$$
        \end{column}
        
    \end{columns}
\end{frame}

\end{document}